\chapter{木管与铜管：用颜色说话，用重音表态}
\label{ch:winds-brass}

\classquote{木管不是“更轻的弦乐”，铜管也不是“更响的弦乐”。它们一个负责颜色和线条，一个负责声明和能量；用错角色，再好的音色也像贴上去的。}{}{}

\section{两组乐器，两种任务}
\label{sec:winds-jobs}

\term{木管组}{woodwind-section}包括长笛、双簧管、单簧管、巴松等；\term{铜管组}{brass-section}包括小号、圆号、长号、大号等。它们最常见的分工是：

\begin{itemize}
\item \textbf{木管}：给旋律换颜色、给和声加线条、在主唱留白处回答；它的优势是灵活、细腻、音色差异大。
\item \textbf{铜管}：在段落转折、副歌进入、结尾收束时“表态”；它的优势是穿透力、厚度和仪式感。
\item \textbf{两者都不该常驻}：木管铺满整首会变成杂音，铜管铺满整首会变成吼叫；它们更适合“进出”，而不是“从头响到尾”。
\end{itemize}

写这两组乐器，先问三个问题：这段需要颜色还是需要能量？它应该在第几拍进入？它和主唱是不是抢同一个音区？答清楚再选乐器。

\section{木管：四种颜色，四条线}
\label{sec:woodwinds}

木管的每件乐器音色差别很大，不能像弦乐那样“同族叠加”：

\begin{itemize}
\item \textbf{长笛}：音区高、线条轻，适合快速跑动、高音副旋律和明亮的色彩；不要让它长期待在中低音区。
\item \textbf{双簧管}：音色集中、带鼻音，穿透力强；适合独奏式的长句和“哭泣”色彩，不宜写成密集伴奏。
\item \textbf{单簧管}：中音区灵活、音色圆，适合连接内声部、与长笛做二重奏；低音区要留出空间。
\item \textbf{巴松}：低音木管，负责木管组的“地板”；适合低音线、次旋律和带簧片的颜色，不适合高速跑动。
\end{itemize}

木管写作最重要的是“每条线有身份”：长笛唱高音副旋律，单簧管填中声部，巴松守低音，双簧管只在关键句出现。四件乐器都吹主旋律，听起来不是“厚”，而是“四个版本的主旋律在打架”。

\section{铜管：声明、重音、厚度}
\label{sec:brass}

铜管的用法更像“标点符号”，不是背景和弦：

\begin{itemize}
\item \textbf{小号}：铜管的最高声部，负责声明、开场、副歌第一拍；音色亮，长时间连续旋律会压过主唱。
\item \textbf{圆号}：音色圆、中音区厚，适合内声部、与弦乐大提琴/中提琴叠合；它是最容易“融进去”的铜管。
\item \textbf{长号}：中低音，适合根音、长音和收束；滑音与短句能增加重量，但不要写成第二套贝斯。
\item \textbf{大号}：最低音，守根音；它不适合旋律，也不该在小声段落强行加入。
\end{itemize}

\term{铜管重音}{brass-stab}是铜管最有效的写法：一小节或两拍的短和弦，落在副歌第一拍、段落转折或尾句；它不需要连续演奏，只需要在关键位置“盖章”。铜管重音太多，会变成每小节一次噪音；太少，又听不出铜管的存在。

\section{音区与穿透力}
\label{sec:winds-ranges}

下表是常用音区，按保守可演奏范围写：

\begin{longtable}{p{0.18\textwidth}p{0.20\textwidth}p{0.42\textwidth}}
\toprule
乐器 & 常用音区 & 配器用途 \\
\midrule
\endfirsthead
\toprule
乐器 & 常用音区 & 配器用途 \\
\midrule
\endhead
长笛 & C4--C7 & 高音副旋律、快速跑动、明亮色彩 \\
双簧管 & Bb3--F6 & 独奏长句、穿透性色彩；少写密集和弦 \\
单簧管 & D3--F6 & 中声部连接、与长笛二重奏 \\
巴松 & Bb1--Bb4 & 木管低音、次旋律；不写高速跑动 \\
小号 & E3--Bb5 & 声明、重音、铜管最高声部 \\
圆号 & F3--C6 & 内声部、与弦乐叠合、圆润色彩 \\
长号 & E2--Bb4 & 根音、长音、收束；不写成第二套贝斯 \\
大号 & E1--Bb3 & 最低根音；不做旋律 \\
\bottomrule
\end{longtable}

\term{穿透力}{projection}不是音量：双簧管和小号在同等推子下会比长笛、圆号更“钻”。编排时先用音区和节奏避让，不要等混音阶段才用推子硬压。铜管整体比木管更容易盖住主唱，尤其当小号与主唱落在同一音区时。

\section{组合：同族、异族、与主唱避让}
\label{sec:winds-combination}

\term{同族组合}{family-grouping}是木管与木管、铜管与铜管的搭配；\term{异族组合}{mixed-grouping}是两组混用。常用做法：

\begin{enumerate}
\item \textbf{木管二重奏}：长笛与单簧管一高一中，吹两条独立线；两条线音区错开，不要同音加倍。
\item \textbf{木管四重奏}：长笛、双簧管、单簧管、巴松各一条线；只在副歌或桥段短时间使用，主歌留白。
\item \textbf{铜管和声}：小号、圆号、长号、大号按“高—中—中低—低”排列；不要四件都吹主旋律。
\item \textbf{木管答句}：主唱句子结束后，用两小节木管回答；答句不重复主唱原句，只接住情绪。
\item \textbf{铜管重音}：副歌第一拍或段落转折给一个短和弦；重音后立即退出，把空间还给主唱。
\end{enumerate}

\term{声部加倍}{doubling}要谨慎：把主旋律原样交给小号或长笛，会让主唱的音区和音色被覆盖。需要加倍时，选择低八度、三度/六度，或只在副歌尾句加倍；不要在整首歌里让木管/铜管与主唱同音同节奏。

\begin{figure}[htbp]
  \centering
  \begin{tikzpicture}[font=\small]
    \draw[->] (0,0) -- (9.2,0) node[right] {bars};
    \foreach \x in {1,...,8} {\draw[gray!30] (\x,0.1) -- (\x,3.3); \node[font=\scriptsize] at (\x,0.05) {\x};}
    \fill[yebanblue!15] (0.2,2.35) rectangle (6.2,2.9);
    \node[font=\scriptsize] at (3.2,2.62) {vocal phrase};
    \fill[yebanblue!55] (2.2,1.45) rectangle (3.2,1.9);
    \fill[yebanblue!55] (4.2,1.45) rectangle (5.2,1.9);
    \node[font=\scriptsize] at (2.7,2.12) {WW answer};
    \node[font=\scriptsize] at (4.7,2.12) {WW answer};
    \fill[orange!55] (6.2,0.55) rectangle (7.2,1.0);
    \fill[orange!55] (7.2,0.55) rectangle (8.2,1.0);
    \node[font=\scriptsize] at (6.7,0.77) {brass};
    \node[font=\scriptsize] at (7.7,0.77) {brass};
    \node[font=\scriptsize, anchor=west] at (8.4,1.67) {woodwinds};
    \node[font=\scriptsize, anchor=west] at (8.4,0.77) {brass};
  \end{tikzpicture}
  \caption{八小节内木管答句与铜管重音的进出：主唱在前六小节，木管只在句尾回答，铜管在副歌转折处盖章。}
  \label{fig:winds-brass-entry}
\end{figure}

\section{在夜半里落地}
\label{sec:winds-yeban}

\begin{enumerate}
\item \textbf{一组一轨，不混装}：长笛、双簧管、单簧管、巴松各建一条 MIDI 轨；小号、圆号、长号、大号各建一条。不要把所有木管塞进一条轨，否则无法单独调平衡。
\item \textbf{装 SFZ 音色}：用「文件 → 装载 SFZ 音色库...」给每条轨装对应音色；同一条轨只放一件乐器，避免音色互相打架。
\item \textbf{在钢琴卷帘里写短句}：按 P / F7 打开卷帘，先写木管答句的两小节，再写铜管重音的一小节；不要把整段写成连续和弦。
\item \textbf{用力度泳道做重音}：铜管重音的第一个音给足力度，后面的音迅速收；木管答句句尾渐弱，不要和主唱抢最后一拍。
\item \textbf{用自动化泳道控制进出}：在编曲区展开自动化泳道（A），让木管在主歌退到很低、在副歌答句时抬起；铜管只在整个重音小节打开。
\item \textbf{用声像与推子收尾}：混音台（X / F8）里把木管与铜管放在不同声像位置；先用编排避让，再用推子修平衡，不要反过来。
\item \textbf{用 Solo/Mute 检查}：Solo 主唱，确认木管/铜管没有同音同节奏覆盖；再 Solo 铜管，确认重音没有盖过鼓的转折。
\end{enumerate}

\screenshot{arrangement-winds-brass}{木管与铜管分轨、钢琴卷帘与自动化泳道同屏：木管答句短、铜管重音短，二者进出位置可见}

\section{常见错误与修正}
\label{sec:winds-mistakes}

\begin{longtable}{p{0.34\textwidth}p{0.58\textwidth}}
\toprule
常见错误 & 修正 \\
\midrule
\endfirsthead
\toprule
常见错误 & 修正 \\
\midrule
\endhead
木管从头铺到尾 & 只在主唱留白处回答；主歌留白，副歌短句进入。 \\
铜管每小节都响 & 改成段落转折/副歌第一拍的短重音，重音后退出。 \\
四件木管都吹主旋律 & 分成高音副旋律、中声部、低音与独奏色彩；每条线有身份。 \\
小号与主唱同音区 & 小号换到低八度或只做短重音；主唱音区留空。 \\
铜管四件都吹和弦 & 按高—中—中低—低分工；不要四件都当旋律。 \\
巴松写快速跑动 & 把快速句交给长笛/单簧管；巴松守低音与长句。 \\
木管/铜管同音加倍 & 改成低八度、三度/六度，或只在副歌尾句加倍。 \\
\bottomrule
\end{longtable}

\section{练习}
\label{sec:winds-exercise}

\begin{exercise}
用八小节写一段“主唱＋木管答句＋铜管重音”：
\begin{enumerate}
\item 先写主唱的四小节乐句，句尾留两拍空白；木管只在空白处写两小节答句。
\item 木管用长笛与单簧管两条线，音区错开；不要复制主唱原句。
\item 第 7 小节写铜管重音：小号、圆号、长号、大号各一个音，按高—中—中低—低排列；只响一拍半。
\item 用力度泳道让铜管重音“进得干脆、退得干净”；再用自动化泳道让木管在主歌几乎听不见。
\end{enumerate}
写完后 Solo 主唱：木管和铜管是否只在应该出现的地方出现？如果它们拿掉后音乐更清楚，说明编排已经完成了它们的任务。
\end{exercise}

\section{本章要点}
\label{sec:winds-takeaways}

\begin{itemize}
\item 木管负责颜色与线条，铜管负责声明与能量；两组都不该常驻。
\item 木管每件乐器音色差异大，要分线写，不要同族无脑叠加。
\item 铜管最有效的是短重音，不是连续和弦。
\item 音区与节奏先避让主唱，混音推子只做最后修补。
\item 同族组合求融合，异族组合求对比；加倍要换音区或音程。
\item 在夜半里一组一轨、SFZ 分装、力度/自动化控制进出，Solo/Mute 验收。
\end{itemize}

\begin{helpmapping}
\begin{longtable}{p{0.45\textwidth}p{0.45\textwidth}}
\toprule
小节 & 帮助锚点 \\
\midrule
两组乐器，两种任务 & \texttt{arrangement-winds-jobs} \\
木管：四种颜色，四条线 & \texttt{arrangement-woodwinds} \\
铜管：声明、重音、厚度 & \texttt{arrangement-brass} \\
音区与穿透力 & \texttt{arrangement-winds-ranges} \\
组合 & \texttt{arrangement-winds-combination} \\
在夜半里落地 & \texttt{arrangement-winds-yeban} \\
常见错误与修正 & \texttt{arrangement-winds-mistakes} \\
练习 & \texttt{arrangement-winds-exercise} \\
\bottomrule
\end{longtable}
\end{helpmapping}
